% Generated draft: requires mathematical and source review.
\documentclass[11pt]{article}
\usepackage[margin=25mm]{geometry}
\usepackage{fontspec}
\setmainfont{Latin Modern Roman}
\usepackage{amsmath,amssymb,mathtools,unicode-math}
\setmathfont{Latin Modern Math}
\usepackage{enumitem,needspace,xurl,hyperref}
\setlength{\parindent}{0pt}
\setlength{\parskip}{5pt}
\setlength{\emergencystretch}{3em}
\newcommand{\N}{\mathbb{N}}
\newcommand{\Z}{\mathbb{Z}}
\newcommand{\Q}{\mathbb{Q}}
\newcommand{\R}{\mathbb{R}}
\newcommand{\C}{\mathbb{C}}
\newcommand{\F}{\mathbb{F}}
\newcommand{\A}{\mathbb{A}}
\newcommand{\PP}{\mathbb{P}}
\newcommand{\E}{\mathbb{E}}
\newcommand{\eps}{\varepsilon}
\newcommand{\dd}{\,\mathrm{d}}
\begin{document}
\section*{UnsolvedMath ID: 30004568}
\textbf{Source problem code:} catalogue entry
\subsection{Definitions and mathematical statement}
The supplied record concerns elastic deformation models without higher-order regularizing terms, a nonlocal penalty for self-contact, and a limiting problem in which self-contact is imposed as a constraint.

For explanatory notation only, let \(X\) denote a space of deformations, let \(P_\eta\) denote a family of nonlocal self-contact penalty energies, and let \(\eta\) approach its limiting value. The source does not specify \(X\), the parameter regime, or the formula for \(P_\eta\). Where a first variation exists, its action on a permissible variation \(h\) at a deformation \(u\) is written
\[
\delta P_\eta(u)[h]
 = \left.\frac{\mathrm d}{\mathrm dt}P_\eta(u+t h)\right|_{t=0},
\]
provided the path and derivative are defined. This notation does not impose a particular differentiability class or identify the first variation with a force without a sign convention.

The source asks the following two questions:
\begin{enumerate}
\item For elastic models without higher-order regularizing terms, can convergence results analogous to the previously considered results be established for the nonlocal self-contact penalty?
\item Along the relevant convergent penalized deformations, do the functional derivatives of the penalty energy converge to the forces caused by self-contact in the constrained limit problem?
\end{enumerate}

These questions cannot be recovered as a fully quantified mathematical conjecture from this record alone. In particular, the antecedent of ``analogous convergence results'' is absent. The elastic energy, reference domain, admissible deformation class, self-contact constraint, penalty scaling, boundary conditions, solution concept, convergence topologies, and mathematical representation of the limiting contact forces are unspecified. The notation above is therefore schematic and does not select these missing hypotheses. Neither convergence of forces nor any particular variational convergence theorem is asserted.
\subsection{Short English statement}
For elastic models without higher-order regularization, the problem asks whether nonlocal self-contact penalties have suitable convergence properties and whether their functional derivatives converge to contact forces in the constrained limit. The supplied record leaves the model and convergence notions unspecified, so a precise theorem statement requires further source information and mathematical review.
\subsection{Sources}
\begin{itemize}
\item Local UnsolvedMath registry, \texttt{data/unsolvedmath/problems.json}: permanent ID \(30004568\), title ``Convergence of Self-Contact Penalty Forces,'' exact source problem code \texttt{catalogue entry}. The supplied record labels the statement \texttt{exact}; its verification note says that the statement is internally readable and complete and was retained byte-for-byte. This does not establish that the missing mathematical context has been recovered.
\item Supplied report attribution: \emph{Mini-Workshop: Nonlocal Analysis and the Geometry of Embeddings} (2021), \emph{Oberwolfach Reports}, pp.\ 1823--1856, DOI \texttt{10.4171/owr/2020/37}. \url{https://doi.org/10.4171/owr/2020/37}. The record separately gives proposed year 2020 and the tag \texttt{owr-year:2020}; no individual proposer is supplied.
\item Supplied literature source: \emph{Surface penalization of self-interpenetration}. \url{https://www.sciencedirect.com/science/article/pii/S0307904X23002731}. No authors, theorem numbers, or detailed bibliographic information are supplied for this item.
\item Supplied literature source: ``OWR 2020/37 report.'' \url{https://doi.org/10.4171/owr/2020/37}.
\item Literature assessment dated \textbf{2026-08-22}: the supplied curation labels the problem ``partially solved'' and states that Gamma-convergence and penalization results exist in restricted linear or regular settings, while convergence of forces and a general low-order model without higher-order terms remain open. This is supplied dated curation derived from the OWR material, not an independently verified literature assessment. Its time-sensitive status and cited claims require independent verification.
\end{itemize}
\end{document}
